\section{Capítulo~\ref{sec:muscles}. A saída para os músculos}

A construção da saída (unidades motoras, margem de segurança da sinapse, ativação por tamanho, par de músculos, laço de estiramento e geradores de ritmo) foi medida diretamente; a condição de estabilidade do laço com atraso é um teorema; o modelo interno do corpo é uma hipótese bem fundamentada; os números das figuras são cálculos pelo modelo do livro.

{\footnotesize
\setlength{\tabcolsep}{3pt}\renewcommand{\arraystretch}{1.15}
\begin{longtable}{>{\raggedright}p{0.5cm}>{\raggedright}p{4.0cm}>{\raggedright}p{5.6cm}>{\raggedright}p{2.3cm}>{\raggedright\arraybackslash}p{1.7cm}}
\toprule
N.º & Proposição & Experimento e o que foi medido & Fonte & Status \\
\midrule\endfirsthead
\toprule
N.º & Proposição & Experimento e o que foi medido & Fonte & Status \\
\midrule\endhead
1 & Unidade motora: um neurônio \nt{ACET} controla um grupo de fibras, de algumas centenas a um par de milhares; nos músculos de ajuste fino as unidades são menores & Músculos humanos, contagem de axônios motores e de fibras musculares: em média cerca de $340$ fibras por neurônio no músculo interósseo da mão, cerca de $410$ no braquiorradial, cerca de $1930$ no gastrocnêmio medial. O capítulo não dá um número exato para as menores unidades. & \cite{Feinstein1955mu} & medido \\
2 & A sinapse no músculo libera várias vezes mais neurotransmissor do que o necessário para a fibra responder & Revisão de medições em sinapses neuromusculares: a margem de segurança (razão entre o neurotransmissor liberado e o limiar) em mamíferos adultos costuma ser $3$--$5$; em humanos, a margem se apoia mais nas dobras da membrana receptora do que no tamanho da liberação. & \cite{WoodSlater2001} & medido \\
3 & O abalo~\eqref{eq:twitch} sobe num tempo $T_c$ da ordem de $50\ms$ & Unidades motoras isoladas da mão humana em esforço voluntário, força promediada pelos disparos da unidade: tempo de subida de $30$--$100\ms$, abaixo de $70\ms$ em $80\,\%$ das unidades. A função $(t/T_c)e^{1-t/T_c}$ é uma aproximação de um modelo de conjunto de unidades. & \cite{MilnerBrown1973rec, Fuglevand1993} & medido \\
4 & Soma de abalos~\eqref{eq:forcesum}: com $5$, $15$ e $40$ disparos/s a força é $0.2$--$1.1F_1$, $1.8$--$2.2F_1$ e cerca de $5.4F_1$ (fig.~\ref{fig:tetanus}) & Cálculo com \texttt{models/\allowbreak muscle\_\allowbreak models.py}, $T_c=50\ms$: $0.21$--$1.10$, $1.76$--$2.19$ e $5.28$--$5.49\,F_1$ nos últimos $0.3\,\text{s}$. A soma linear é uma simplificação: o modelo não tem a saturação do músculo real. & \texttt{models/\allowbreak muscle\_\allowbreak models.py} & modelo \\
5 & As unidades são ativadas por tamanho; as mais fracas são cem vezes mais fracas que as mais fortes & Raízes ventrais do gato: com entrada reflexa crescente, disparam primeiro os neurônios com impulsos pequenos na raiz, isto é, os pequenos. Músculo interósseo da mão humana: força de abalo das unidades de $0.1$ a $10$~g, crescendo quase linearmente com a força em que a unidade é ativada. & \cite{Henneman1957, MilnerBrown1973rec} & medido \\
6 & O passo de força é proporcional à força, $\Delta F/F\approx q-1$~\eqref{eq:sizeprinciple}: ponto flutuante & Deduzido no capítulo a partir da progressão geométrica das forças. Concorda com o experimento: em esforços grandes, o mesmo incremento de força ativa muito menos unidades novas. & dedução no capítulo; \cite{MilnerBrown1973rec} & teoria \\
7 & A dispersão da força cresce proporcionalmente à própria força & Esforço isométrico voluntário do músculo do polegar: o desvio padrão da força cresce linearmente com a força média; com o mesmo esforço provocado por estimulação elétrica, a dispersão é constante. Modelo de conjunto: o crescimento linear resulta da ativação das unidades por ordem de força. & \cite{Jones2002sdn} & medido \\
8 & A suavidade dos movimentos é consequência do ruído dependente do sinal & Modelo: a trajetória que minimiza a dispersão do ponto final com esse ruído reproduz as trajetórias medidas do olho e do braço. Não há experimento direto que distinga essa causa das outras. & \cite{HarrisWolpert1998} & hipótese \\
9 & A diferença das forças do par de músculos dá o torque, a soma dá a rigidez~\eqref{eq:antagonists} & Teoria do controle de impedância por co-contração. Braço humano: a rigidez de cada articulação, medida com pequenos empurrões, é aproximadamente proporcional ao torque nela; diferentes proporções de co-contração mudam a forma e a direção da elipse de rigidez da mão. & \cite{Hogan1984, GomiOsu1998stiff} & medido \\
10 & O sensor de estiramento excita o seu neurônio \nt{ACET} e, por um intermediário inibitório, desliga o antagonista (fig.~\ref{fig:antagonists}) & Medula espinhal do gato: um único intermediário, excitado pelas fibras dos sensores de estiramento, produz potenciais inibitórios monossinápticos nos neurônios \nt{ACET} do músculo antagonista, atraso sináptico de $0.28$--$0.42\ms$. O neurotransmissor do intermediário (glicina) não foi determinado nesse experimento. & \cite{Jankowska1972Ia} & medido \\
11 & Atraso do laço: espinhal $25$--$30\ms$, pelo córtex uma vez e meia a três vezes mais, pelo olho $100$--$200\ms$ & Braço humano, empurrão na articulação: resposta muscular curta após $20$--$45\ms$, longa após $45$--$105\ms$, voluntária após $120$--$180\ms$. Deslocamento da posição visível do dedo durante o movimento: correção após $160\ms$ em média. & \cite{Pruszynski2008rapid, SaundersKnill2003vis} & medido \\
12 & $\dd x/\dd t=-Gx(t-d)$ é estável com $Gd<\pi/2$~\eqref{eq:delaystable}; no limite, a frequência é $1/(4d)$ & Teorema sobre as raízes de uma equação com retardo. Verificação com \texttt{models/\allowbreak muscle\_\allowbreak models.py}, $d=30\ms$: em $3\,\text{s}$, amplitude $3\cdot10^{-6}$ com $G=40$, $0.13$ com $G=50$, $38$ com $G=55$ por segundo (limite $52$). & \cite{Hayes1950} & teoria \\
13 & Tremor fisiológico de $8$--$12$~Hz; o laço de estiramento é uma das suas fontes & Revisão de medições: um tremor fino na faixa de cerca de $10$~Hz existe em pessoas saudáveis; a sua origem é mista (oscilações centrais, propriedades das descargas das unidades, ressonância mecânica e ressonância do laço reflexo), e em condições diferentes predomina uma ou outra. & \cite{McAuleyMarsden2000} & hipótese \\
14 & O cérebro prevê o resultado do comando por um modelo interno do corpo & Uma pessoa segura um peso em pinça e move o braço: com carga inercial, viscosa e elástica, a força de preensão muda ao mesmo tempo que a carga, e não com o atraso do laço, ou seja, a carga é prevista. Uma revisão reúne esses experimentos; não há observação direta do próprio modelo nos neurônios. & \cite{FlanaganWing1997grip, WolpertGhahramani2000} & hipótese \\
15 & O ritmo de andar, respirar e mastigar é gerado por redes locais com entrada constante & Revisão de experimentos: o sistema nervoso isolado (medula, gânglios) produz um padrão motor rítmico sem entrada rítmica e sem realimentação dos músculos. & \cite{MarderBucher2001} & medido \\
16 & A inibição mútua com fadiga~\eqref{eq:halfcentre} dá um ritmo com período de $0.84\,\text{s}\approx2\tau_a$ (fig.~\ref{fig:halfcentre}) & Teorema: uma rede com inibição mútua e adaptação sem estado estável, com entrada constante, necessariamente oscila. Cálculo com \texttt{models/\allowbreak muscle\_\allowbreak models.py} ($I=1$, $\beta=2$, $b=2$, $\tau=20\ms$, $\tau_a=0.4\,\text{s}$): período de $0.84\,\text{s}$. Que os geradores reais sejam construídos exatamente assim não foi mostrado. & \cite{Matsuoka1985osc} & modelo \\
\bottomrule
\end{longtable}}
